\section{Ventajas y costos: por qué el Transformer tiene éxito y por qué sigue siendo difícil}

El éxito del Transformer suele atribuirse a las dependencias de largo alcance y al entrenamiento en paralelo, pero ambas ventajas tienen condiciones. La full attention ofrece rutas de información cortas, pero genera $n^2$ score; el cálculo matricial se adapta bien a GPU/TPU, pero exige lotes grandes, memoria y conjuntos de datos de gran escala; un sesgo inductivo local más débil aumenta la generalidad, pero también puede incrementar la cantidad de muestras necesarias.

\subsection{Complejidad y límites del sistema}

En una sola capa de full self-attention, el cálculo de los score es de aproximadamente $O(n^2d)$ y la score memory, de aproximadamente $O(n^2)$; una sola capa de RNN suele costar alrededor de $O(nd^2)$ y requiere $O(n)$ sequential steps. Cuál de las dos es más rápida depende de $n$, $d$, el lote, el kernel, el memory bandwidth y el hardware. Una tabla de complejidad no equivale directamente a la wall-clock latency.

\lecturefigure{slide-15-good-bad.jpg}{El modelado de largo alcance y el paralelismo del Transformer vienen acompañados de una attention cuadrática, de la demanda de datos y de un sesgo local más débil.}{Video oficial, 16:10.}

\begin{knowledgebox}{Lectura de la figura: las ventajas y las desventajas provienen del mismo diseño}
La interacción directa entre cualquier par de token permite capturar dependencias de largo alcance, pero también produce una dense pairwise matrix; no depender de la recurrence permite entrenar en paralelo, pero también reduce el sesgo de orden incorporado; un block uniforme facilita el escalamiento, pero puede requerir más datos para aprender la estructura del dominio. Al leer la figura, no hay que tratar las dos columnas como listas independientes, sino buscar las dos caras de una misma decisión de diseño.
\end{knowledgebox}

\begin{warningbox}{FlashAttention no vuelve lineal la complejidad matemática}
La fused/tiled attention puede reducir las lecturas y escrituras en \term{HBM} (High Bandwidth Memory, memoria de video de alto ancho de banda de la GPU) y la materialización de matrices intermedias, lo que mejora notablemente la velocidad y el uso de memoria en la práctica; sin embargo, el cálculo pairwise de la full attention estándar sigue siendo cuadrático. La algorithmic sparsity, el low-rank, los state-space o la recurrence son enfoques distintos; la optimización del sistema y la complejidad algorítmica deben describirse por separado.
\end{warningbox}

\begin{knowledgebox}{Términos clave: inductive bias y sample efficiency}
El \term{inductive bias} (sesgo inductivo) es la estructura que el modelo prefiere antes de ver los datos; por ejemplo, el compartir pesos local e invariante a traslaciones de la convolution, la recursión temporal de la RNN o la direccionalidad de la causal mask. Un sesgo menor ofrece mayor generalidad, pero el modelo puede necesitar más datos para aprender regularidades que podrían haberse codificado a mano. La \term{sample efficiency} describe la cantidad de datos necesaria para alcanzar el desempeño objetivo, no el número de muestras procesadas por segundo.
\end{knowledgebox}

\subsection{Resumen de la sección}

Las ventajas y los costos del Transformer tienen el mismo origen: la conexión global dinámica, el paralelismo matricial y los supuestos débiles sobre el dominio amplían la escalabilidad, pero también traen un costo cuadrático, demanda de datos y complejidad del sistema. El juicio de ingeniería debe considerar al mismo tiempo el algoritmo y el hardware.
